\documentclass[10pt,a4paper]{article}
\usepackage{amsmath,amssymb,amsthm,graphicx}
\title{A logarithmic-gap dichotomy for randomized nonadaptive reductions to Kolmogorov complexity}
\author{Michael Saks, Rahul Santhanam}
\date{doi:10.4230/LIPIcs.CCC.2022.29; CC BY4.0}
\begin{document}
\maketitle
\noindent\textit{Editorial note.} Original authors: Michael Saks and Rahul Santhanam. Official CCC2022 article29 (2022-07-11), DOI10.4230/LIPIcs.CCC.2022.29. Creative Commons Attribution4.0.\par\medskip
\noindent\textit{Editorial note.} License: https://creativecommons.org/licenses/by/4.0/\par\medskip
\noindent\textit{Editorial note.} Faithful editable TeX transcription of formal Theorem3 with complete author Sections4--6 profile-oracle, AM protocol and dichotomy proofs, exact required definitions and complete bibliography. Original published mathematical expressions retained, without correction or new proofs. Theorem30 and reverse Corollary32 remain grouped with the main result.\par\medskip
\noindent\textit{Allowed standard prerequisites (editorial).}\begin{itemize}
\item Entropy estimation and approximate preimage lower-bound AM protocols. Exact cited author Lemmas21[24] and22[25] retained.
\item Hirahara generator distinguishability reduction. Exact cited author Theorem31[28,Theorem6.3] retained, including honesty clarification.
\end{itemize}\medskip
\paragraph{Author descriptions and Kolmogorov complexity (Section 2.4).}
Throughout we assume a fixed pairing function $\langle\cdot\rangle$ and a standard encoding of Turing machines $M$ by binary strings, with $\widetilde M$ denoting the encoding of $M$.
We fix a universal Turing machine $U$ that for any Turing machine $M$ and input $x$ takes as input $\langle\widetilde M,x\rangle$ and outputs $M(x)$. The Kolmogorov complexity $K(q)$ of $q$ (with respect to $U$), denoted $K(q)$, is the minimum length of $y$ such that $U(y)=q$. We define the Kolmogorov complexity of a machine $M$ to be $K(\widetilde M)$.
\paragraph{Proposition 11 (Symmetry of information).}
Let $x,y$ be strings. Then $K(x;y)=K(x)+K(y\mid x)+O(\log(K(x,y)))$.
\paragraph{Author random-string promise problems (Section 2.5).}
In the $R_K$ problem, the instance is a string $x$, and the question is whether $K(x)\geq|x|/2$. Given a function $g:\mathbb N\to\mathbb N$, the $R_K$ problem with gap $g$ is the promise problem whose NO instances are strings $x$ with $K(x)\geq|x|/2$ and whose YES instances are strings $x$ with $K(x)<|x|/2-g(|x|)$. Given a function $\alpha:\mathbb N\to\mathbb R$ such that $\alpha(n)\geq1$ for all $n$, the $R_K$ problem with multiplicative gap $\alpha$ is the promise problem whose NO instances are strings $x$ with $K(x)\geq|x|/2$ and whose YES instances are strings $x$ such that $K(x)<|x|/(2\alpha(|x|))$.

\paragraph{Author oracle-machine conventions (Sections 2.2--2.3).}
We assume without loss of generality that a polynomial-time bounded randomized Turing Machine $M$ comes equipped with a poly-time computable function $r^M:\{0,1\}^*\to\mathbb N$ where $r_n^M$ is a polynomial upper bound on the number of random bits generated by the machine on input $x$. On input $x$ the TM starts by generating an auxiliary random string $\rho$ of length $r_x^M$ which is written on the random tape, and the program operates deterministicially on the pair $x,\rho$.
Normally, an oracle is a function $O$ from a query set $Q$ to an answer set $A$. We also need a more general notion of a context-sensitive oracle in which the response depends on the query $q$ and the context (machine $M$ and input $x$ that is calling the oracle). We write $O_x^M(q)$ for the output of oracle $O$ to query $q$.
Unless otherwise specified all oracles in this paper have query set $Q=\{0,1\}^*$ and answer set $A=\mathbb N$.
An OTM $M$ instantiated by oracle $O$, denoted $M^O$ is the program that uses oracle $O$ to answer any queries. If $O$ is computable, then $M^O$ is computably instantiated, and can be identified with an oracle-free Turing machine, even in the case that $O$ is a context-sensitive oracle.
The randomized OTMs we consider are decision machines which means that they always halt and output either 1 (identified with accept) or 0 (identified with reject). The language $L(M^O)$ recognized by an (possibly randomized) instantiated decision OTM $M^O$ is the set of $x$ for which the probability that $M(x)$ accepts is greater than $1/2$.
Let $\epsilon:\mathbb N\to[0,1]$. We say that $L$ is accepted by $M^O$ with advantage $\epsilon$ if every accepted string $x$ is accepted with probability at least $(1+\epsilon(|x|))/2$ and every rejected string is accepted with probability at most $(1-\epsilon(|x|))/2$. Note that computing with advantage $\epsilon=1$ means that $M^O(x)$ always outputs $L(x)$.
An algorithm that accepts $L$ with advantage $\epsilon$ can be converted to an algorithm with advantage $(1-\epsilon')$ by repeating the algorithm $O(\log(1/\epsilon')+(1/\epsilon))^2$ times and taking majority vote.
We restrict attention to OTMs that are nonadaptive. A nonadaptive OTM comes equipped with:
\begin{itemize}
\item A computable function $k^M:\{0,1\}^*\to\mathbb N$. For $x\in\{0,1\}^*$, $k_x^M$ is the number of oracle queries that $M$ makes on input $x$. For convenience (and with no significant loss of generality) we assume that $k_n^M$ is a power of 2.
\item A computable function $q^M$ that takes as input a string $x$ and an integer $i\in\{1,\ldots,k_{|x|}^M\}$, and (if the OTM is randomized) a random string $\rho$ of length $r_{|x|}^M$, and outputs an element of $\{0,1\}^*$. $q_x^M(i,\rho)$ is the $i$th query asked by $M$ on input $x$ when the random string is $\rho$.
\item A computable function $\operatorname{out}^M$ that takes as input $x$ and the sequence of query answers $(a_1,\ldots,a_{k_x^M})$ and (if the OTM is randomized) the random string $\rho$ and outputs either 0 or 1. $\operatorname{out}_x^M(a_1,\ldots,a_{k_x^M},\rho)$ is the output of the algorithm on input $x$, random string $\rho$ when the query answers are $r_1,\ldots,r_{k_x^M}$.
\end{itemize}
In the above definition the nonadaptive OTM $M$ is deterministic if $r_x^M$ is identically 0. Thus the function $q_x^M$ depends only on the query index $i$.
$Q_x^M$ denotes the set of all $q\in\{0,1\}^*$ such that on input $x$, $M$ asks query $q$ with positive probability. $Q_x^M$ is a finite set of size at most $2^{r_x^M}k_x^M$ since there are $2^{r_x^M}$ choices for $\rho$ and each results in $k_x^M$ queries.
We say that $M$ is $\alpha$-honest, for $\alpha:\mathbb N\to\mathbb N$, provided that on any input $x$, all oracle queries have length at least $\alpha(|x|)$. If $\alpha(n)=n^\tau$ for some constant $\tau>0$ we say that $M$ is polynomially honest. Conventionally, a polynomially honest oracle algorithm is referred to simply as honest, but it will be useful for us to distinguish different kinds of honesty.
If $M$ is $\alpha$-honest for some $\alpha$ that tends to $\infty$ we say that $M$ is weakly honest. We say that $M$ is fixed query length if the length of each query of $M$ on input $x$ depends only on $|x|$.
\paragraph{2.3. Approximate Oracles and Robust Oracle Algorithms.}
Recall that in this paper we consider oracles that take as input a string and output an integer. We now introduce two related concepts for an oracle (possibly context-sensitive) $O$:
\begin{itemize}
\item Approximation of $O$ by another oracle $O'$ that is ``suitably close''.
\item An oracle algorithm $M$ instantiated by $O$ that is robust with respect to replacing $O$ by any $O'$ that is close enough to $O$.
\end{itemize}
The closeness of two oracles is measured by a nonnegative valued closeness function $\beta$ with two arguments, a string $q$ and natural number $n$. The value of $\beta$ at these arguments is written $\beta_n(q)$. We say that $O$ and $P$ are $\beta$-close with respect to $M$ provided that for all inputs $x$ to $M$, and queries $q\in Q_x^M$, $|O_x^M(q)-P_x^M(q)|\leq\beta_{|x|}(q)$.
Let $M$ be a randomized oracle algorithm, $O$ be an oracle and $L$ be the language accepted by $M^O$. Let $\beta$ be a closeness function and let $\epsilon:\mathbb N\to[0,1]$. We say that $M$ is $(\beta,\epsilon)$-robust with respect to oracle $O$ provided that for every oracle $O'$ that is $\beta$-close to $O$, $M^{O'}$ accepts $L$ with advantage $\epsilon$. We say that $M$ is $\beta$-robust if it is $(\beta,1)$-robust, which means that for every oracle $O'$ that is $\beta$-close to $O$, $M^{O'}(x)$ always outputs the correct answer $L(x)$.
\paragraph{Proposition 7.}
For every computable function $f:\{0,1\}^*\to\{0,1\}^*$ there is a constant $C_f$ such that if $q=f(z)$ then $K(q)\leq K(z)+C_f$.
In particular, for $f$ being the identity function $\mathsf{ID}$ we have
\paragraph{Proposition 8.}
For all strings $q$, $K(q)\leq|q|+C_{\mathsf{ID}}$.
\paragraph{Proposition 9.}
$\sum_{q\in\{0,1\}^*}2^{-K(q)}/K(q)^2\leq2$.
\begin{proof}
Letting $S_k$ be the set of strings of Kolmogorov complexity exactly $k$, we can rewrite the above sum as:
\[\sum_{k\geq1}\sum_{q\in S_k}\frac{2^{-k}}{k^2}.\]
Since $|S_k|\leq2^k$, this is at most $\sum_{k\geq1}1/k^2<2$.
\end{proof}
The canonical order on $\{0,1\}^*$ is the order in which $x$ precedes $y$ if $|x|<|y|$ or $|x|=|y|$ and $x$ is lexicographically less than $y$.
\paragraph{Proposition 10.}
Let $L$ be a computable language that contains infinitely many strings.
\begin{enumerate}
\item Let $\Lambda:\mathbb N\to\mathbb N$ be any computable unbounded function. Let $x_j$ be the first string (according to the canonical order) satisfying $|x_j|\geq\Lambda(j)$ and $x\in L$. There is an integer $C$ (depending on $L$ and $\Lambda$) such that for all integers $j$, $K(x_j)\leq\log(j)+C$.
\item Let $\kappa:\mathbb N\to\mathbb Z$ be any computable nondecreasing function that tends to $\infty$. Let $L$ be a computable language that contains infinitely many strings. Then $L$ contains a string $x$ for which $K(x)<\kappa(|x|)$.
\end{enumerate}
\begin{proof}
Let $M$ be a Turing Machine that computes $L$ and $N$ be a Turing machine that computes $\Lambda$. We can construct a machine $M'$ that on input $j$ outputs $x_j$. $M$ first evaluates $\Lambda(j)$ and then enumerates strings $x$ of length at least $\Lambda(j)$ in the canonical order and stops when it finds the first string in $L$. $M'$ is a program whose length is a constant (depending on $|N|$ and $|M|$) and so the pair $M',j$ has a description of length $\log(j)+C$.
For the second part, given $\kappa$, let $\Lambda$ be the function where $\Lambda(j)$ is the least integer $m$ such that $2\log(j)<\kappa(m)$. The function $\Lambda$ is computable. Applying the first part, there is a constant $C$ (depending on $L$ and $\Lambda$) such that for all $j$, $x_j$ is a string of length at least $\Lambda(j)$ such that $K(x_j)\leq\log(j)+C$. Choose $j=2^C$. Then $x_j$ has length at least $\Lambda(j)$, and $K(x_j)\leq\log(j)+C\leq2\log(j)<\kappa(\Lambda(j))\leq\kappa(|x_j|)$, as required.
\end{proof}

\section*{Formal author main result}
\paragraph{Theorem 3.} If a language $L$ is computable and there is a randomized polynomial-time non-adaptive reduction $F$ from $L$ to $\omega(\log(n))$-additively approximating Kolmogorov complexity such that $F$ is polynomially honest and has inverse polynomial advantage, then $L\in\mathrm{AM}\cap\mathrm{coAM}$. On the other hand, every $L\in\mathrm{NEXP}$ has a randomized polynomial-time non-adaptive reduction making $q$ queries to $O(\log(n))$-additively approximating Kolmogorov complexity that is polynomially honest and has constant advantage.

\paragraph{4. Simulating robust $K$-oracle computation in $\mathsf{AM}\cap\mathsf{coAM}$.}
Throughout this section we make the following assumptions.
\begin{description}
\item[H1] $M$ is a randomized nonadaptive decision OTM.
\item[H2] $M^K$ recognizes the computable language $L$.
\item[H3] $\beta=\beta_n(q)$ is a computable closeness function.
\item[H4] $\epsilon:\mathbb N\to[0,1]$.
\item[H5] $M$ is $(\beta,\epsilon)$-robust with respect to oracle $K$.
\item[H6] $\alpha:\mathbb N\to\mathbb N$ is an unbounded non-decreasing function.
\end{description}
Also, let $k_n^M$ denote the maximum of $k_x^M$ over $x$ of length $n$.
Our overall goal is to prove the following theorem:
\paragraph{Theorem 14.}
Suppose $M,\alpha,\beta,\epsilon$ satisfy hypotheses [H1--H6]. Suppose also that $\beta_n(q)=\log(k_n^M/\epsilon(n))+2\log K(q)+\alpha(n)+5$. Then both $L$ and $\overline L$ have constant round Arthur--Merlin protocols for membership whose running time is polynomial in the running time of $M$ and $1/\epsilon$. In particular if $M$ runs in polynomial time then $L\in\mathsf{AM}\cap\mathsf{coAM}$.
The rather cumbersome condition on $\beta$ arises from the proof. It is more natural if the closeness function depends only on the query $q$ and not on the length of the input $x$ to the calling algorithm.
We prove Theorem 14 in two parts. First, in this section, we construct a parameterized family $\{J[\gamma]:\gamma\in[0,1)\}$ of computable context-sensitive oracles $J$ based on simple combinatorial properties of the machine $M$ and show that there exists an integer $\ell$ such that under [H1]--[H6], $M^{[\ell]J[\gamma]}$ computes $L$ with advantage $\epsilon/2$ for every $\gamma\in[0,1)$.
Then, in the next section, we show that if $L$ satisfies the conclusion of the first part, then it has an efficient $\mathsf{AM}$ protocol and $\mathsf{coAM}$ protocol.
\paragraph{4.1. Some additional preliminaries.}
Let $(M,x)$ be a context. As usual, for a set or quantity $Z$ that depends on the context, we denote this dependence by $Z_x^M$, but we often suppress this dependence and write simply $Z$. For example we write $r$ for $r_x^M$ and $k$ for $k_x^M$.
Let $\Gamma=\Gamma_x^M=\{(\rho,i):(\rho,i)\in\{0,1\}^r\times\{1,\ldots,k\}\}$. Under the assumption stated earlier that $k$ is a power of 2 we have that $\log|\Gamma|=r+\log k$ is an integer, and we can identify $\Gamma$ with the set $\{0,1\}^{r+\log k}$.
It will be convenient to introduce notation for some elementary functions of positive real numbers. For a nonnegative integer $s$:
\begin{itemize}
\item $\phi$ is the function given by $\phi(s)=\lceil\log(|\Gamma|/s)\rceil=\log(|\Gamma|)-\lfloor\log s\rfloor$.
\item $\mu$ is the function given by: $\mu(s)$ is the least nonnegative number $\gamma$ such that $s(1+\gamma)$ is a power of 2. Note that $\mu(s)\in[0,1)$.
\end{itemize}
It is easy to check:
\paragraph{Proposition 15.}
For $\gamma\in[0,1)$ and $s>0$, $\phi((1+\gamma)s)=\phi(s)$ if $\gamma<\mu(s)$ and $\phi((1+\gamma)s)=\phi(s)-1$ if $\gamma\geq\mu(s)$.
The following technical fact will be needed later.
\paragraph{Proposition 16.}
Let $s,t\in\mathbb N$ and $\tau\in[0,1)$ with $t\in[s/(1+\tau),s(1+\tau)]$. Let $\gamma\in[0,1)$ be such that $\gamma\notin[\mu(s)-2\tau,\mu(s)+2\tau]\cup[1-2\tau,1)\cup[0,2\tau)$. Then $\phi((1+\gamma)s))=\phi((1+\gamma)t)$.
\begin{proof}
We prove the contrapositive. Suppose $\phi((1+\gamma)s)\neq\phi((1+\gamma)t)$. Then $\lfloor\log((1+\gamma)s)\rfloor\neq\lfloor\log((1+\gamma)t)\rfloor$. This means that there is an integral power of 2 between them, which means that the interval $[(1+\gamma)s/(1+\tau),(1+\gamma)(1+\tau)s]$ contains an integral power of 2, say $2^w$. Note that $s\leq2^{w+1}<8s$. This implies that $(1+\mu(s))s=2^{w+1}$, or $(1+\mu(s))s=2^w$, or $(1+\mu(s))s=2^{w-1}$.
Case 1. Suppose $(1+\mu(s))s=2^{w+1}$. This means $s>2^w$, and since $[(1+\gamma)s/(1+\tau),(1+\gamma)(1+\tau)s]$ contains $2^w$, we have that $\tau>\gamma$. Hence $\gamma\in[0,2\tau)$.
Case 2. Suppose $(1+\mu(s))s=2^w$. Then $(1+\mu(s))s$ belongs to the interval $[(1+\gamma)s/(1+\tau),(1+\gamma)(1+\tau)s]$ so $1+\mu(s)$ belongs to the interval $[(1+\gamma)/(1+\tau),(1+\gamma)(1+\tau)]\subseteq[1+\gamma-2\tau,1+\gamma+2\tau]$ which implies $\gamma\in[\mu(s)-2\tau,\mu(s)+2\tau]$.
Case 3. Suppose $(1+\mu(s))s=2^{w-1}$. Then $2\leq2(1+\mu(s))=2^w/s\leq(1+\gamma)(1+\tau)\leq1+\gamma+2\tau$ which implies $\gamma\geq1-2\tau$.
\end{proof}
\paragraph{4.2. Replacing $K$ by the oracle $J[\gamma]$.}
As described earlier the context $(M,x)$ determines a function $q_x^M:\Gamma\to\{0,1\}^*$ where $q_x^M(\rho,i)$ is the $i$th query asked by $M$ on input $x$ when the random string is $\rho$.
For a query $q$, we define $S(q)=S_x^M(q)$ to be the set $\{(\rho,i)\in\Gamma:q_x^M(\rho,i)=q\}$.
Let $s(q)=|S(q)|$ and let $\pi(q)=\pi_x^M(q)=s_x^M(q)/|\Gamma|$. This is the probability, with respect to $(\rho,i)$ chosen uniformly from $\Gamma$ that $q(\rho,i)=q$. Note that for the function $\phi$ defined earlier we have $\phi(s(q))=\lceil\log(1/\pi(q))\rceil$.
For each $\gamma\in[0,1)$, define the function $J[\gamma]$ on queries by:
\[J[\gamma](q)=J[\gamma]_x^M(q)=\phi((1+\gamma)s(q)).\]
If $\gamma=0$ we often write $J(q)$ for $J[\gamma](q)$. From the definition we observe that for all $\gamma\in[0,1)$:
\[J[\gamma](q)\in[\log(1/\pi(q))-1,\log(1/\pi(q))+1].\tag{2}\]
We view $J[\gamma]$ as a context-sensitive oracle and we will see below that $J[\gamma]_x^M$ can be used in place of $K$ in a suitably robust oracle algorithm. The following Lemma relates $J[\gamma]$ to $K$:
\paragraph{Lemma 17.}
Let $M$ be a randomized oracle algorithm and $x$ a string. Recall that $Q_x^M$ is the set of queries $q$ such that $M$ on input $x$ asks the query $q$ with positive probability.
\begin{enumerate}
\item There are constants $C_0,C_1$ such that for any string $q\in Q_x^M$, $K(q)-J[\gamma]_x^M(q)\leq C_0K(x)+C_1$.
\item For $u>0$ let $P_x^M(u)$ be the set of strings in $Q_x^M$ such that $J[\gamma]_x^M(q)-K(q)\geq u+2\log K(q)$. Then $\pi_x^M(P_x^M(u))\leq2^{3-u}$.
\end{enumerate}
\begin{proof}
Fix $M$ and $x$. We omit the superscript $M$ and subscript $x$ from $J[\gamma],\pi,P$ and $Q$.
For the first part, order the elements of $Q$ lexicographically by the triple $(1/\pi(q),|q|,q)$. For a query $q$, let $p(q)$ be the number of $q'\in Q$ that precede $q$ in this order. Given the Turing machine $M$ and input $x$ we can construct a Turing machine of size $O(K(x))+O(1)$ that takes as input an integer $h\in\{0,\ldots,|Q|-1\}$ and outputs $q$ such that $p(q)=h$.
Therefore $K(q)\leq O(K(x))+\log p(q)+O(1)$. Since there are at most $1/\pi(q)$ members of $Q$ that precede $q$ in the order, we have $p(q)\leq(1/\pi(q))$ and therefore $\log p(q)\leq\log(1/\pi(q))\leq J[\gamma](q)+1$. Therefore $K(q)-J[\gamma](q)\leq O(K(x))+O(1)$, as required.
For the second part, (2) implies $\pi(q)\leq2^{1-J[\gamma](q)}$. For $q\in P(u)$ we therefore have $\pi(q)\leq2^{1-K(q)-u}/K(q)^2$. Summing over all $q\in P(u)$ this is at most $2^{1-u}\sum_{q\in P(u)}2^{-K(q)}/K(q)^2$ which is at most $2^{2-u}$ by Proposition 9.
\end{proof}
\paragraph{4.3. Replacing a $K$ oracle by a $J[\gamma]$ oracle.}
In this section we prove:
\paragraph{Theorem 18.}
Assume hypotheses [H1]--[H6] and that $\beta_{|x|}(q)=\log(k_x^M/\epsilon(|x|))+2\log K(q)+\alpha(|x|)+5$. Then there is an integer $\ell$ such that for all $\gamma\in[0,1]$, $M^{[\ell]J[\gamma]}(x)$ computes $L$ with advantage $\epsilon/2$, where $M^{[\ell]}$ is the OTM that outputs $L(x)$ on any input $x$ of length less than $\ell$ and otherwise executes $M$ on $x$. Letting $N=M^{[\ell]}$, there is an OTM $N$ such that $N^{J[\gamma]}(x)$ computes $L$ with advantage $\epsilon/2$.
\begin{proof}
Fix $M$. Say that $x$ is bad for $\gamma$ if $M^{J[\gamma]}(x)$ outputs $L(x)$ with probability less than $\frac12(1+\epsilon(|x|)/2)$ and let $B[\gamma]$ be the set of strings that are bad for $\gamma$. Let $B=\bigcup_{\gamma\in[0,1)}B[\gamma]$.
If $B$ is finite, let $\ell$ be the length of the largest string in $B$. Then $M^{[\ell]}$ satisfies the desired conclusion. So we assume that $B$ is infinite and derive a contradiction.
Note that $x\in B$ holds if and only if one of the infinitely many conditions $x\in B[\gamma]$ holds. The following Proposition reduces this to a finite set of conditions.
\paragraph{Proposition 19.}
Let $x\in\{0,1\}^*$ and let $\Gamma_x=\{\mu(s(q)):q\in Q_x^M\}$ (where $Q_x^M$ is the (finite) set of queries that are asked with nonzero probability by $M$ on input $x$ and $\mu$ is the function defined in the definition of $J[\gamma]$). Then $x\in B$ if and only if $x\in B[\gamma]$ for some $\gamma\in\Gamma_x$.
\begin{proof}
For each $\gamma\in[0,1)$ let $\gamma^-=\max\{\mu(s(q)):q\in Q_x^M,\mu(s(q))\leq\gamma\}$. By the definition of $J[\gamma]$, we have $J[\gamma](q)=J[\gamma^-](q)$ for all $q\in Q$. Therefore $M^{J[\gamma]}(x)$ behaves identically to $M^{J[\gamma^-]}(x)$ and so $x\in B$ if and only if there exists $\gamma\in\Gamma_x$ such that $x\in B[\gamma]$.
\end{proof}
Let us define $\gamma_x$ to be the least $\gamma\in\Gamma_x$ for which $x\in B[\gamma_x]$.
The set $\Gamma_x$ is computable from $x$, and for each $\gamma\in\Gamma_x$ we can computably determine the probability that $M^{J[\gamma]}(x)$ agrees with $L(x)$. Therefore there is an (instantiated) program of size $|M|+O(1)$ that on input $x$, tests whether $x\in B$.
Let $\beta^*(n)$ be the minimum over all $q$ and all $n'\geq n$ of $\beta_{n'}(q)$. To obtain the desired contradiction we will prove:
\paragraph{Lemma 20.}
For any string $x$ there is an oracle $K[x]$ satisfying:
\begin{itemize}
\item $K[x]$ is $\beta$-close to $K$ and therefore the probability that $M^{K[x]}$ on input $x$ differs from $M^K$ on input $x$ is at most $(1-\epsilon(|x|))/2$.
\item If $K(x)<(\beta^*(|x|)-C_1)/C_0$ (where $C_0,C_1$ are given in Lemma 17) then the probability that $M^{K[x]}$ on input $x$ differs from $M^J$ on input $x$ is at most $\epsilon(|x|)/4$.
\end{itemize}
The lemma implies any string $x$ satisfying $K(x)<(\beta^*(|x|)-C_1)/C_0$ is not in $B$, since the probability that $M^J$ on input $x$ disagrees with $M^K$ on input $x$ can be upper bounded by $(1-\epsilon(|x|)/2+\epsilon(|x|)/4\leq(1-\epsilon(|x|)/2)/2$. But by applying the second part of Proposition 10 to the language $B$ and $\kappa(n)=(\beta^*(n)-C_1)/C_0$ there is a string $x\in B$ such that $K(x)<(\beta(|x|)-C_1)/C_0$, which yields the desired contradiction. Note that $\kappa$ is an unbounded non-decreasing function as required for the application of Proposition 10 by the definition of $\beta^*$ and the fact that $\alpha$ is an unbounded non-decreasing function.
So it remains to prove Lemma 20. Given $x$, we define the oracle $K[x]$ as follows. Recall that $Q_x^M$ is the set of queries $q$ that are made by $M$ with non-zero probability when the input is $x(w)$. Let $Q_{\mathrm{close}}$ be the set of $q\in Q_x^M$ for which $|J[\gamma]_x^M(q)-K(q)|\leq\beta_{|x|}(q)$, let $Q_>$ be the set of $q\in Q_x^M$ for which $J[\gamma_x]_x^M(q)>K(q)+\beta_{|x|}(q)$ and let $Q_<$ be the set of $q\in Q_x^M$ for which $J[\gamma_x](q)<K(q)-\beta_{|x|}(q)$. Define the oracle $K[x]$ as follows: $K[x](q)=J[\gamma_x](q)$ if $q\in Q_{\mathrm{close}}$ and $K[x](q)=K(q)$ otherwise.
By definition $K[x]$ is $\beta$-close to $K$ and therefore by hypothesis $M^{K[x]}$ computes $L$ with advantage $\epsilon$. In particular $M^{K[x]}$ on input $x$ differs from $L(x)$ with probability at most $(1-\epsilon(|x|))/2$. This completes the first part of the lemma.
For the second part of the lemma, assume that $x$ is such that $K(x)<(\beta(|x|)-C_1)/C_0$ (where $C_0,C_1$ are given in Lemma 17).
If all of the queries asked by $M$ on input $x$ are in $Q_{\mathrm{close}}$ then all query answers during the computation $M^{J[\gamma]}(x)$ are identical to those during the computation $M^{K[x]}(x)$, and therefore the output is the same.
So it now suffices to prove an upper bound of $\epsilon(x)/4$ on the probability that at least one of the selected queries on input $x$ belongs to $Q_<\cup Q_>$.
From the hypothesis, we have $K(x)\leq(\beta(|x|)-C_1)/C_0$. We now show that this implies that $Q_<=\varnothing$. By the first part of Lemma 17, for all queries $q\in Q_x^M$ we have $K(q)-J_x^M(q)\leq C_0K(x)+C_1$ which is at most $\beta(|x|)$ and so $q\notin Q_<$.
Hence we only need to upper bound the probability that, on input $x$, at least one query belongs to $Q_>$. For brevity let $k=k_x^M$. For $i\in\{1,\ldots,k\}$ let $p_i$ be the probability that the $i$th query is in $Q_>$. Let $\sigma=\sum_ip_i$, it suffices to show that $\sigma\leq\epsilon(|x|)/4$. Note that $\sigma/k$ is the chance that a random query $i$ (with respect to the random bits and the choice of $i\in\{1,\ldots,k\}$) is in $Q_>$. By the second part of Lemma 17, with $u=\beta(|x|)-2\log(K(q))$, $p\leq2^{2+2\log(K(q))-\beta(|x|)}$. Using the hypothesis on $\beta(|x|)$ this is at most $\epsilon(|x|)/4k$ and therefore $\sigma\leq\epsilon(|x|)/4$, as required.
\end{proof}

\section*{The Arthur--Merlin protocol (author Section 5)}
The starting point for this section is the conclusion of Theorem 18:
\paragraph{Hypothesis AM$(N)$.} $N$ is an oracle machine such such that for all $\gamma\in[0,1)$, $N^{J[\gamma]}(x)$ computes $L$ with advantage $\epsilon/2$.
We will show that this implies that there is an Arthur--Merlin protocol for $L$ whose running time is in the running time of $N$.
We make use of two known AM protocols. The input to both protocols includes a function $f:\{0,1\}^w\longrightarrow\{0,1\}^*$ given by a circuit $C_f$. The function $f$ induces a probability distribution $\pi_f$ on $\{0,1\}^*$, with $\pi_f(q)$ equal to $|f^{-1}(q)|/2^2$, which is the probability that a uniformly chosen element of $\{0,1\}^w$ maps to $q$. Also $\zeta,\delta>0$ are input parameters for the protocol.
\paragraph{Lemma 21 (Entropy Estimation Protocol, [24]).} There is a constant round AM-protocol that takes as input a circuit $C_f$ as above, an integer $h$ and an error parameter $\delta$, runs in time $|C_f|\operatorname{poly}(\log(1/\delta))$, and outputs accept or reject with the following conditions:
\begin{enumerate}\item If $|h-H(\pi_f)|\le1$ then Merlin has a strategy that causes Arthur to reject with probability at most $\delta$.
\item If $|h-H(\pi_f)|\ge2$ then the probabilty that the verification is accepted is at most $\delta$.\end{enumerate}
\paragraph{Lemma 22 (Lower Bound protocol, [25]).} There is a constant round protocol with input $C_f$, input sequences $q_1,\ldots,q_m\in(\{0,1\}^*)^m$ and $(s_1,\ldots,s_m)\in\mathbb N^m$, and error parameters $(\zeta,\delta)$ both in $(0,1)$ that runs in time $C_f\operatorname{poly}(m,1/\zeta,\log(1/\delta))$, and outputs either accept or reject according to the following:
\begin{enumerate}\item If $s_i\ge|f^{-1}(q_i)|$ for every $i$, then Merlin has a strategy such that the protocol rejects with probability at most $\delta$.
\item If $s_i\le|f^{-1}(q_i)|/(1+\zeta)$ for some $i$ then for any strategy of Merlin the probabilty that the verification is accepted is at most $\delta$.\end{enumerate}
We fix some notation for this section: $x$ denotes the string whose membership in $L$ is to be determined. $r=r_x^N$ is the length of the random string generated by $N$ on input $x$. $k=k_x^N$ is the number of queries asked by $N$ on input $x$. We assume (with no significant loss of generality) that $\log k$ is an integer. For $\rho\in\{0,1\}^r$, $q_1(\rho),\ldots,q_k(\rho)$ denotes the queries asked by $N$ on input $x$. $f$ denotes the function on $\{0,1\}^{r+\log k}$ that maps $(\rho,i)$ to $q_i(\rho)$. Let $\pi=\pi_f$ be the probability distribution on $\{0,1\}^*$ induced by $f$, which is given by $\pi_f=|f^{-1}(q)|/2^{r+\log k}$. $Q=Q_x^N$ is the image of the function $f$, which is the set of queries that are asked with positive probability. For $\rho\in\{0,1\}^r$, $a_1[\gamma](\rho),\ldots,a_k[\gamma](\rho)$ is the sequence of answers from the oracle $J[\gamma]$ to the queries $q_1(\rho),\ldots,q_k(\rho)$. When the random string is $\rho$, the output of the algorithm is $\operatorname{out}(a_1[\gamma](\rho),\ldots,a_k[\gamma](\rho);\rho)$. We denote this by $z[\gamma](\rho)$. $\epsilon=\epsilon_{|x|}^N$.
\paragraph{Proposition 23.} Let $\pi$ be a probability distribution over $\{0,1\}^w$.
\begin{enumerate}\item If $q$ is selected according to $\pi$, $\mathbb E[\log(1/\pi(q))]=H(\pi)$ where $H(\pi)$ is the binary entropy of $\pi$.
\item Suppose $q_1,\ldots,q_t$ are selected independently according to $\pi$. Then
\[\Pr\left[\left|\frac1t\sum_{i=1}^t\log(1/\pi(q_i))-H(\pi)\right|>\zeta\right]\le\frac{w^2}{t\zeta^2}.\]\end{enumerate}
\begin{proof}The first part follows directly from the definition of $H(\pi)$. For the second part, let $Z_i=\log(1/\pi(q_i))$ and let $Z=\frac1t\sum Z_i$. Note that $\mathbb E[Z]=\frac1t\sum\mathbb E[Z_i]=H(\pi)$ by the first part. Also, note that since $Z$ is an average of $t$ i.i.d. random variables, it's variance is $\frac1t\operatorname{Var}(Z_i)\le w^2/t$ since a random variable taking values in $[0,w]$ has variance at most $w^2$. By Chebyshev's inequality, for any random variable $X$, the probability that $|Z-\mathbb E[Z]|>\zeta$ is at most $\operatorname{Var}(Z)/\zeta^2\le w^2/(t\zeta^2)$.\end{proof}
In describing the Arthur--Merlin protocol we refer to the actions of ``Honest Merlin'' which is the ideal behavior for Merlin (which Merlin may not follow). The protocol makes use of several parameters, whose values are determined by the analysis. For reference we list these parameters here with their values:
\begin{itemize}\item $t$ is the sample complexity, and it set to $t=32000\max(r,1)^2k^2/\epsilon(|x|)^2$.
\item $\tau$ is the safety parameter, which is set to $\tau=\epsilon(|x|)/(16(4k+2))$.
\item $\nu$ is the lower bound parameter, which is set to $\nu=\epsilon(|x|)^2/(20000k^2)$.
\item $D$ is the precision parameter and is a positive integer. We choose $D=\lceil16(4k+2)/\epsilon(|x|)\rceil$.\end{itemize}
It is straightforward to verify that these parameter values satisfy the following conditions:
\paragraph{Proposition 24.}
\begin{align*}
\tau&\le\frac{\epsilon(|x|)}{16(4k+2)},& \nu&\le\tau\epsilon(|x|)/(128k),\\
t&\ge\frac{256k}{\epsilon(|x|)}\log\frac{16}{\epsilon(|x|)},& t&\ge\frac{16(r+\log k)^2}{\epsilon(|x|)},\\
t&\ge\frac{320k}{\epsilon(|x|)\tau},&D&\ge\frac1\tau,\\
D&\ge\frac{32(k+1)}{\epsilon(|x|)}.&&
\end{align*}
The inequalities stated in this Proposition are needed in the analysis of the protocol. The values of the parameters were chosen to ensure these conditions. We are now ready to state the protocol:
\paragraph{Step 1.} Merlin provides an integer $h$. Arthur and Merlin perform the Entropy estimation protocol of Lemma 21 for the probability distribution $\pi$ and input $h$ with error parameter $\delta=\epsilon(|x|)/16$. (If Merlin is honest, then Merlin sets $h$ so that $|h-H(\pi)|\le1$, and adopts the strategy from the first part of Lemma 21 that makes the rejection probability at most $\epsilon(|x|)/16$.)
\paragraph{Step 2.} Random strings $\rho_1,\ldots,\rho_t$, each of length $r$, are generated using public coins. Let $\chi$ be the $t\times k$ matrix with $\chi_{i,j}=f(\rho_i,j)$, the $j$th query asked when the random string is $\rho_i$.
\paragraph{Step 3.} Merlin provides a $t$ by $k$ positive integer matrix $B$ (Honest Merlin sets $B$ to be equal to the matrix $B^*$ whose $i,j$ entry is $|f^{-1}(\chi_{i,j})|$.)
\paragraph{Notational Remark.} It is convenient here to define some matrices and vectors that are determined from $B$. For $\gamma\in[0,1)$ the $t\times k$ matrix $A[\gamma]$ is given $A[\gamma]_{i,j}=\phi((1+\gamma)B_{i,j})$, and the vector $z[\gamma]\in\{0,1\}^t$ is given by $z[\gamma]_i=\operatorname{out}(A[\gamma]_{i,1},\ldots,A[\gamma]_{i,k})$. The matrix $A^*[\gamma]$ and vector $z^*[\gamma]$ are obtained analogously with $B$ replaced by $B^*$. Note that row $i$ of $A^*[\gamma]$ is exactly the sequence of query answers provided by oracle $J[\gamma]$ when the random string is $\rho_i$ and $z^*[\gamma]_i$ is the output of $N^{J[\gamma]}(x)$ for random string is $\rho_i$.
\paragraph{Step 4.} For each row index $i$, Arthur randomly chooses $\hat j(i)\in\{1,\ldots,k\}$. Define
\[\sigma(B)=\frac1t\sum_i(r-\log(B_{i,\hat j(i)})).\]
Arthur computes $\sigma(B)$ approximately to determine an integer $H'$ that belongs to the interval $[\sigma(B)-1,\sigma(B)+1]$. Arthur rejects this step if $|H'-h|>3$ and accepts this step otherwise.
\paragraph{Step 5.} Arthur--Merlin perform the lower bound protocol of Lemma 22 for the above-mentioned function $f$ that maps $(\rho,i)$ to $q_i(\rho)$, with error parameters $(\zeta,\delta)=(\nu,\epsilon(|x|)/16)$. The input $(q_1,\ldots,q_m),(s_1,\ldots,s_m)$ to the protocol is the sequence of $tk$ entries of the matrix $\chi$ and the corresponding sequence of entries of the matrix $B$. (If Merlin is Honest then since $B=B^*$, $s_i=|f^{-1}(q_i)|$ for every $i$ and Merlin follows the strategy from the first part of Lemma 22 so that the probability of rejection is at most $\epsilon(|x|)/16$.)
\paragraph{Step 6.} Arthur selects a uniformly random $\hat\lambda\in\{1,\ldots,D\}$ and a uniformly random $\hat i\in\{1,\ldots,t\}$. $\hat z$ is defined to be $z[\lambda/D]_{\hat i}$. For future reference we also define $\hat z^*$ to be $\hat z$ if $B$ is replaced by $B^*$.
\paragraph{Step 7.} If any of the steps 1,4 or 5 result in reject, then the protocol outputs failure. Otherwise the output is $\hat z$.
This protocol outputs either failure or 0 or 1. The main theorem of this subsection is:
\paragraph{Theorem 25.} Under the hypotheses AM$(N)$ given earlier, the above protocol satisfies:
\begin{enumerate}\item If Merlin behaves honestly, then the protocol outputs $L(x)$ with probability at least $\frac12+\epsilon(|x|)/16$.
\item For any behavior of Merlin, the probability that the protocol outputs failure or outputs $L(x)$ is at least $\frac12+\epsilon(|x|)/16$.\end{enumerate}
\paragraph{Corollary 26.} Both $L$ and $\overline L$ can be recognized by AM protocols that run in time polynomial in the $1/\epsilon$ and the running time of $N$.
\begin{proof}The protocol obtained by changing a failure output to 0 (resp. 1) is an AM protocol (resp. coAM protocol) for $L$.\end{proof}
\paragraph{Proof of Theorem 25.} We first note the following:
\paragraph{Proposition 27.} The value $\hat z^*$ determined in Step 6 in the case that $B=B^*$ agrees with $L(x)$ with probability at least $\frac12+\epsilon(|x|)/4$.
\begin{proof}The value $\hat z^*$ determined in Step 6 is equal to the output of $N^{J[\lambda/D]}$ on input $x$ when the random string is $\rho_{\hat i}$. Since $\rho_{\hat i}$ is a uniformly random string, the hypothesis on $N$ implies that this is equal to $L(x)$ with probability at least $\frac12+\epsilon(|x|)/4$.\end{proof}
We now prove the first part of the Theorem. For $s\in\{1,4,5\}$ define the event $R_s$ to be the event that the test in Step $s$ outputs reject. Assume Merlin behaves honestly, and therefore $B=B^*$. By the above proposition, the probability that the output of the protocol is $L(x)$ is at least:
\begin{align*}
&\frac12+\frac{\epsilon(|x|)}4-\Pr[R_1\lor R_4\lor R_5]\\
&\qquad\ge\frac12+\frac{\epsilon(|x|)}4-\Pr[R_1]-\Pr[R_4]-\Pr[R_5].
\end{align*}
As noted in the description of Steps 1 and 5, $\Pr[R_1]$ and $\Pr[R_5]$ are each at most $\epsilon(|x|)/16$. So it suffices to show that $\Pr[R_4]\le\epsilon(|x|)/16$, which we now do. Step 4 is rejected only if $|H'-h|>3$. We have:
\[|H'-h|\le|H'-\sigma(B^*)|+|\sigma(B^*)-H(\pi)|+|H(\pi)-h|.\]
$|H'-\sigma(B^*)|\le1$ by the choice of $H'$ and the fact that $B=B^*$, and $|H(\pi)-h|\le1$ since Merlin is assumed to be honest. Therefore $\Pr[R_4]\le\Pr[|\sigma(B^*)-H(\pi)|>1$. So it suffices to prove:
\paragraph{Proposition 28.} $\Pr[|\sigma(B^*)-H(\pi)|>1]\le\epsilon(|x|)/16$.
\begin{proof}By the definition of $B^*$, $B^*_{i,\hat j}=2^{r+\log k}\pi(\chi_{i,\hat j(i)})$ and so $r+\log k-\log(B^*_{i,\hat j(i)})=\log(1/\pi(\chi_{i,\hat j(i)}))$ and so $\sigma(B^*)=\frac1t\sum_{i=1}^t\log(1/\pi(\chi_{i,\hat j(i)}))$. Since $\chi_{1,\hat j(1)},\ldots,\chi_{t,\hat j(t)}$ is a sequence of independent samples from the distribution $\pi$, Proposition 23 implies
\[\Pr\left[\left|\frac1t\sum_i\log(1/\pi(\chi_{i,\hat j(i)}))-H(\pi)\right|>1\right]<\frac{(r+\log k)^2}t\le\epsilon(|x|)/16,\]
where the last inequality used Proposition 24.\end{proof}
This completes the proof of the first part of the theorem. We now turn to the proof of the second part. We consider an arbitrary strategy for Merlin. We break up into cases according to properties of the matrix $B$. We divide into three cases.
\paragraph{Case 1.} There is at least one entry $(i,j)$ with $B_{i,j}<B^*_{i,j}/(1+\nu)$. By the property of the Lower bound protocol performed in Step 5, this will reject with probability at least $1-\epsilon(|x|)/16\ge\frac12)+\epsilon(|x|)/4$.
\paragraph{Case 2.} $B_{i,j}\ge B^*_{i,j}/(1+\nu)$ for all entries $(i,j)$ and there are more than $\epsilon(|x|)t/16$ entries with $B_{i,j}>(1+\tau)B^*_{i,j}$. We claim that Step 4 will reject with probability at least $1-\epsilon(|x|)/8\ge1/2+\epsilon/4$. For this we give a lower bound on the probability $H'-H(\pi)$ exceeds 3. We have
\begin{align*}H'-H(\pi)&=\sigma(B)-\sigma(B^*)-(\Sigma(B)-H')-(H(\pi)-\sigma(B^*))\\
&\ge\sigma(B)-\sigma(B^*)-|\sigma(B)-H'|-|H(\pi)-\sigma(B^*)|.\end{align*}
By the choice of $H'$, $|\sigma(B)-H'|\le1$. Therefore:
\begin{align*}\Pr[H'-H(\pi)>3]&\ge\Pr[(\sigma(B)-\sigma(B^*)\ge5)\land(|\sigma(B^*)-H(\pi)|<1)]\\
&\ge1-\Pr[\sigma(B)<\sigma(B^*)+5]\\&\qquad-\Pr[|\sigma(B^*)-H(\pi)|\ge1].\end{align*}
The desired lower bound will follow by showing that each of the two probabilities on the right is at most $\epsilon(|x|)/16$. Proposition 28 shows this for the second probability. So we now bound the first probability. For $i\in\{1,\ldots,t\}$ define:
\begin{itemize}\item $V_i=\{j\in\{1,\ldots,k\}:\log B_{i,j}\ge(1+\tau)B^*_{i,j}$. (Note that $\sum_i V_i\ge\epsilon(|x|)t/16$ by the case assumption.)
\item $Y_i$ is the 0-1 indicator of the event that $\hat j(i)\in V_i$, and $Y=\sum_iY_i$.\end{itemize}
We now make two claims. Claim 1. $\sigma(B)-\sigma(B^*)\ge\tau Y-\nu t$. Claim 2. $\Pr[Y<\frac12\frac{\epsilon(|x|)t}{16k}]\le\epsilon(|x|)/16$.
If $Y\ge\frac12\frac{\epsilon(|x|)t}{16k}$, then by the conditions on $\tau$ and $\nu$ given by Proposition 24 we have that $\tau Y-\nu t\ge5$ and so
\[\Pr[\sigma(B)-\sigma(B^*)<5]\le\Pr\left[Y<\frac12\frac{\epsilon(|x|)t}{16k}\right]\le\epsilon(|x|)/16.\]
To prove the first claim, note that for $i$ such that $Y_i=1$, we have
\[\log(B_{i,\hat j(i)})\ge\log((1+\tau)B^*_{i,\hat j(i)}\ge\log(B^*_{i,\hat j(i)})+\tau\]
(since $\log(1+\tau)\ge\tau$ for $\tau\in[0,1]$). For $i$ such that $Y_i=0$, we have $B_{i,\hat j(i)}\ge B_{i,\hat j(i)}/(1+\nu)$ which implies $\log B_{i,\hat j(i)}\ge\log B_{i,\hat j(i)}-2\nu$ since $\log(1+\nu)\le2\nu$). Therefore $\sigma(B)-\sigma(B^*)\ge\tau Y-\nu t$.
For the second claim, note that that $\mathbb E[Y]=\sum_i\mathbb E[Y_i]=\sum_i|V_i|/k\ge\epsilon(|x|)t/(16k)$, by the case assumption. The ``multiplicative Chernoff bound'' implies that if $Y$ is a sum of independent 0-1 random variables, then $\Pr[Y\le\frac12\mathbb E[Y]]\le e^{-\mathbb E[Y]/8}$. Applying this here we obtain $\Pr[Y\le\epsilon(|x|)t/(32k)\le e^{-\epsilon(|x|)t/(256k)}\le\epsilon(|x|)/16$ by Proposition 24.
\paragraph{Case 3.} $B_{i,j}\ge B^*_{i,j}/(1+\nu)$ for all entries $(i,j)$ and there are at most $\epsilon(|x|)t/16$ entries with $B_{i,j}>(1+\tau)B^*_{i,j}$. By Proposition 27, $\hat z^*$ is equal to $L(x)$ with probability at least $\frac12+\epsilon(|x|)/4$. We show that under the case assumption the probability that $\hat z\ne\hat z^*$ is at most $\epsilon(|x|)/8$ which will show that $\hat z$ is equal to $L(x)$ with probability at least $\frac12+\epsilon(|x|)/8$.
We say that a row index $i\in\{1,\ldots,t\}$ is unsafe if there is at least one $j$ such that $B_{i,j}>(1+\tau)B^*_{i,j}$ and is safe otherwise. We identify the following ``bad'' events:
\begin{itemize}\item[B1] The selected row index $\hat i$ is unsafe.
\item[B2] The selected $\hat\lambda$ satisfies
\[\hat\lambda/D\in[0,2\tau)\cup[1-2\tau,1)\cup\bigcup_{j=1}^k[\mu(B^*_{\hat i,j})-2\tau,\mu(B^*_{\hat i,j})+2\tau].\]\end{itemize}
We claim that if neither bad event happens then $\hat z=\hat z^*$. Since [B1] doesn't happen, $\hat i$ is safe. Together with the case assumption and the fact that $\nu\le\tau$ this implies that for each $j\in[k]$, $B_{\hat i,j}\in[B^*_{\hat i,j}/(1+\tau),B^*_{\hat i,j}(1+\tau)]$. Since [B2] doesn't happen, then for all $j\in\{1,\ldots,k\}$, the hypotheses of Proposition 16 hold for $(s,t)=(B^*_{\hat i,j},B_{\hat i,j})$ from which we conclude that for all $j\in\{1,\ldots,k\}$
\begin{align*}A[\lambda/D]_{\hat i,j}&=\phi((1+\lambda/D)B_{\hat i,j}))\\&=\phi((1+\lambda/D)B^*_{\hat i,j}))=A^*[\lambda/D]_{\hat i,j},\end{align*}
which implies that $\hat z[\lambda/D]=\hat z^*[\lambda/D]$. So now it will be enough to show that $\Pr[\mathrm{B1}]$ and $\Pr[\mathrm{B2}]$ are both at most $\epsilon(|x|)/16$. $\Pr[\mathrm{B1}]$ is equal to the fraction of unsafe rows, which is at most $\epsilon(|x|)/16$ by the case assumption.
Since $\lambda$ is chosen uniformly from $\{1,\ldots,D\}$, $\Pr[\mathrm{B2}]$ is the fraction of integers in $\{1,\ldots,D\}$ that belong to the set
\begin{multline*}[0,2D\tau)\cup[D(1-2\tau),D)\\\cup\bigcup_{j=1}^k[D(\mu(B^*_{\hat i,j})-2\tau),D(\mu(B^*_{\hat i,j})+2\tau)].\end{multline*}
Now we use the fact that the number of integers in an interval $[a,b)$ is less than $b-a+1$ and an interval $[a,b]$ is at most $b-a+1$. Therefore the number of integers in $[0,2D\tau)\cup[D(1-2\tau),D)$ is less than $4\tau D+2$ and the number of integers in each set of the big union is at most $4\tau D+1$. Summing up over all intervals the number of integers in the union is at most $(4k+4)\tau D+k+1$. Since $\tau D\le1/4$ by Proposition 24, this is at most $2k+2$. So the probability that $\lambda$ is in this set is at most $(2k+2)/D$ which is at most $\epsilon(|x|)/16$ by Proposition 24.\qed

\section*{A Dichotomy (author Section 6)}
We first make Theorem 14 more concrete by identifying some interesting parameter settings. One natural setting is where the number of queries is polynomial and the advantage is inverse polynomial.
\paragraph{Theorem 29.} Suppose $L$ is computable and there is a randomized polynomial-time non-adaptive reduction $F$ from $L$ to $\omega(\log(|q|))$-additively approximating Kolmogorov complexity such that $F$ is polynomially honest and has inverse polynomial advantage. Then $L\in\mathrm{AM}\cap\mathrm{coAM}$.
\begin{proof}We apply Theorem 14 with $k_n^M=\operatorname{poly}(n)$, $\epsilon(n)=1/\operatorname{poly}(n)$, and $\alpha(n)=O(\log(n))$. Hence we have $\beta_n(q)=O(\log(n))$. Since $F$ is polynomially honest, we have that $\beta_n(q)=O(\log(q))$ for every $q\in Q_x^M$ for $x$ of length $n$. Hence by the conclusion of Theorem 14, we have $L\in\mathrm{AM}\cap\mathrm{coAM}$.\end{proof}
Another natural setting is where the number of queries and the advantage are both constant. In this case, we can get a conclusion from a smaller upper bound on the approximation gap.
\paragraph{Theorem 30.} Suppose $L$ is computable and there is a computable function $\eta=\omega(1)$ such that there is a randomized polynomial-time non-adaptive reduction $F$ from $L$ to $(2\log(K(q))+\eta(|q|))$-additively approximating Kolmogorov complexity such that $F$ makes $O(1)$ queries, is weakly honest and has advantage $\Omega(1)$. Then $L\in\mathrm{AM}\cap\mathrm{coAM}$.
\begin{proof}Let $\lambda:\mathbb N\to\mathbb N$ be an unbounded function such that $F$ is $\lambda$-honest. When applying Theorem 14, we choose $k_n^M=O(1)$ and $\epsilon(n)=\Omega(1)$. We also choose $\alpha$ appropriately so that $\beta_n(q)\ge2\log(K(q))+\eta(|q|)$---this is possible since $F$ is weakly honest, by ensuring $\alpha(\lambda^{-1})=o(\eta)$. Once more, applying Theorem 14 gives us the desired condlusion.\end{proof}
Theorem 29 supplies the first part of our dichotomy result. We require the following result of Hirahara [28] for the second part of our dichotomy. This is Theorem 6.3 in [28], where it is stated without mentioning the honesty property of the reduction. The proof of Theorem 6.3 in [28] does yield a polynomially honest reduction.
\paragraph{Theorem 31 ([28]).} For any constant $c\in\mathbb N$, there exists an $\mathrm{EXP}^{\mathrm{NP}}$ computable function $G=\{G_n\}$, where $G_n$ maps $n-c\log(n)$ bits to $n$ bits, such that for any $L\in\mathrm{NEXP}$, there is a randomized polynomial-time polynomially honest non-adaptive reduction from $L$ to distinguishing $G$ from uniform.
\paragraph{Corollary 32.} For every constant $c\in\mathbb N$ and any $L\in\mathrm{NEXP}$, there is a randomized polynomial-time non-adaptive reduction $F$ from $L$ to $c\log(|q|)$-additively approximating Kolmogorov complexity such that $F$ is polynomially honest and has constant advantage.
\begin{proof}Let $c\in\mathbb N$, and let $K'$ be any oracle that $c\log(|q|)$-additively approximates Kolmogorov complexity. Wlog we assume $c\ge3$---note that establishing the conclusion for a constant $c$ also establishes it for any constant $c'\le c$. Applying Theorem 31, we have that there is a computable function $G=\{G_n\}$, where $G_n$ maps $n-3c\log(n)$ bits to $n$ bits such that for any $L\in\mathrm{NEXP}$, there is a randomized polynomial-time polynomially honest non-adaptive reduction from $L$ to $D$, where $D$ is any language distinguishing $G$ from the uniform distribution. We define $D(q)=1$ if $K'(q)<|q|-1.5c\log(|q|)$. Note that any output of the generator $G$ can be described with $\log(n)+n-3c\log(n)+O(1)$ bits, and hence has Kolmogorov complexity $<n-2.5c\log(n)$ for large enough $n$. This implies that for every $n$-bit output $y$ of the generator $G_n$, $K'(y)<n-1.5c\log(n)$. On the other hand, with probability at least $1-1/n^{c/2}$ over $y\sim U_n$, we have that $K(y)\ge n-0.5c\log(n)$, and hence $K'(y)\ge n-1.5c\log(n)$. Thus $K'$ does distinguish the output of $G_n$ from uniform. Since $D'$ can be computed with a single deterministic query to $K'$ of the same length as the input, we have that for any $L\in\mathrm{NEXP}$, there is a randomized polynomial-time polynomially honest non-adaptive reduction from $L$ to $K'$. Since this holds for any $K'$ that $c\log(|q|)$-additively approximates Kolmogorov complexity, the result follows.\end{proof}
Theorem 29 together with Corollary 32 establishes the dichotomy result in Theorem 3.

\begin{thebibliography}{99}
\bibitem[1]{published-1} Adi Akavia, Oded Goldreich, Shafi Goldwasser, and Dana Moshkovitz. On basing one-way functions on np-hardness. In Proceedings of the 38th Annual ACM Symposium on Theory of Computing, Seattle, WA, USA, May 21-23, 2006, pages 701–710, 2006.
\bibitem[2]{published-2} Eric Allender. The complexity of complexity. In Computability and Complexity - Essays Dedicated to Rodney G. Downey on the Occasion of His 60th Birthday, pages 79–94, 2017.
\bibitem[3]{published-3} Eric Allender. The new complexity landscape around circuit minimization. In Proceedings of 14th International Conference on Language and Automata Theory and Applications, volume 12038 of Lecture Notes in Computer Science, pages 3–16. Springer, 2020.
\bibitem[4]{published-4} Eric Allender, Harry Buhrman, Luke Friedman, and Bruno Loff. Reductions to the set of random strings: The resource-bounded case. Log. Methods Comput. Sci., 10(3), 2014.
\bibitem[5]{published-5} Eric Allender, Harry Buhrman, and Michal Koucký. What can be efficiently reduced to the kolmogorov-random strings? Ann. Pure Appl. Log., 138(1-3):2–19, 2006.
\bibitem[6]{published-6} Eric Allender, Harry Buhrman, Michal Koucký, Dieter van Melkebeek, and Detlef Ronneburger. Power from random strings. SIAM J. Comput., 35(6):1467–1493, 2006.
\bibitem[7]{published-7} Eric Allender, Harry Buhrman, Michal Koucký, Dieter van Melkebeek, and Detlef Ronneburger. Power from random strings. SIAM J. Comput., 35(6):1467–1493, 2006.
\bibitem[8]{published-8} Eric Allender and Bireswar Das. Zero knowledge and circuit minimization. In Symposium on Mathematical Foundations of Computer Science (MFCS), pages 25–32, 2014. 7 Technically speaking, it might be unclear if Lemma 35 applies here, as the sampled strings do not in general have length m, but the proof of Lemma 35 continues to work as long as the sampled strings have length m$\Omega$(1)
\bibitem[9]{published-9} Eric Allender, George Davie, Luke Friedman, Samuel Hopkins, and Iddo Tzameret. Kolmogorov complexity, circuits, and the strength of formal theories of arithmetic. Chic. J. Theor. Comput. Sci., 2013, 2013.
\bibitem[10]{published-10} Eric Allender, Luke Friedman, and William I. Gasarch. Limits on the computational power of random strings. Inf. Comput., 222:80–92, 2013.
\bibitem[11]{published-11} Eric Allender, Joshua A. Grochow, and Cristopher Moore. Graph isomorphism and circuit size. CoRR, abs/1511.08189, 2015. arXiv:1511.08189.
\bibitem[12]{published-12} Eric Allender, Lisa Hellerstein, Paul McCabe, Toniann Pitassi, and Michael E. Saks. Minimizing disjunctive normal form formulas and AC0 circuits given a truth table. SIAM J. Comput., 38(1):63–84, 2008.
\bibitem[13]{published-13} Eric Allender and Shuichi Hirahara. New insights on the (non-)hardness of circuit minimization and related problems. In International Symposium on Mathematical Foundations of Computer Science (MFCS), pages 54:1–54:14, 2017.
\bibitem[14]{published-14} Eric Allender, Dhiraj Holden, and Valentine Kabanets. The minimum oracle circuit size problem. In International Symposium on Theoretical Aspects of Computer Science (STACS), pages 21–33, 2015.
\bibitem[15]{published-15} Eric Allender, Michal Koucký, Detlef Ronneburger, and Sambuddha Roy. The pervasive reach of resource-bounded Kolmogorov complexity in computational complexity theory. J. Comput. Syst. Sci., 77(1):14–40, 2011.
\bibitem[16]{published-16} Luis Filipe Coelho Antunes and Lance Fortnow. Worst-case running times for average-case algorithms. In Proceedings of the 24th Annual IEEE Conference on Computational Complexity, CCC 2009, Paris, France, 15-18 July 2009, pages 298–303. IEEE Computer Society, 2009.
\bibitem[17]{published-17} Benny Applebaum, Boaz Barak, and David Xiao. On basing lower-bounds for learning on worst-case assumptions. In Proceedings of 49th Annual IEEE Symposium on Foundations of Computer Science, pages 211–220, 2008.
\bibitem[18]{published-18} Sanjeev Arora and Boaz Barak. Computational Complexity: A Modern Approach. Cambridge University Press, 1st edition, 2009.
\bibitem[19]{published-19} Sergei Artemenko, Russell Impagliazzo, Valentine Kabanets, and Ronen Shaltiel. Pseudorandomness when the odds are against you. In Proceedings of the 31st Conference on Computational Complexity, volume 50 of LIPIcs, pages 9:1–9:35. Schloss Dagstuhl - Leibniz- Zentrum für Informatik, 2016.
\bibitem[20]{published-20} Andrej Bogdanov and Luca Trevisan. On worst-case to average-case reductions for NP problems. SIAM J. Comput., 36(4):1119–1159, 2006.
\bibitem[21]{published-21} Harry Buhrman, Lance Fortnow, Michal Koucký, and Bruno Loff. Derandomizing from random strings. In Proceedings of the 25th Annual Conference on Computational Complexity, CCC ’10, 2010.
\bibitem[22]{published-22} Marco L. Carmosino, Russell Impagliazzo, Valentine Kabanets, and Antonina Kolokolova. Agnostic learning from tolerant natural proofs. In Approximation, Randomization, and Combinatorial Optimization, pages 35:1–35:19, 2017.
\bibitem[23]{published-23} Joan Feigenbaum and Lance Fortnow. Random-self-reducibility of complete sets. SIAM J. Comput., 22(5):994–1005, 1993.
\bibitem[24]{published-24} Oded Goldreich and Salil P. Vadhan. On the complexity of computational problems regarding distributions. In Oded Goldreich, editor, Studies in Complexity and Cryptography, volume 6650, pages 390–405. Springer, 2011.
\bibitem[25]{published-25} Shafi Goldwasser and Michael Sipser. Private coins versus public coins in interactive proof systems. In Juris Hartmanis, editor, Proceedings of the 18th Annual ACM Symposium on Theory of Computing, pages 59–68. ACM, 1986.
\bibitem[26]{published-26} Dan Gutfreund and Salil P. Vadhan. Limitations of hardness vs. randomness under uniform reductions. In Proceedings of 12th International Workshop on Randomness and Computation, volume 5171 of Lecture Notes in Computer Science, pages 469–482. Springer, 2008.
\bibitem[27]{published-27} Shuichi Hirahara. Non-black-box worst-case to average-case reductions within NP. In Mikkel Thorup, editor, 59th IEEE Annual Symposium on Foundations of Computer Science, FOCS 2018, Paris, France, October 7-9, 2018, pages 247–258. IEEE Computer Society, 2018.
\bibitem[28]{published-28} Shuichi Hirahara. Unexpected hardness results for kolmogorov complexity under uniform reductions. In Proceedings of 52nd Annual Symposium on Theory of Computing, pages 1038–1051, 2020.
\bibitem[29]{published-29} Shuichi Hirahara, Igor Carboni Oliveira, and Rahul Santhanam. Np-hardness of minimum circuit size problem for OR-AND-MOD circuits. In 33rd Computational Complexity Conference, CCC 2018, pages 5:1–5:31, 2018.
\bibitem[30]{published-30} Shuichi Hirahara and Osamu Watanabe. Limits of minimum circuit size problem as oracle. In Conference on Computational Complexity (CCC), pages 18:1–18:20, 2016.
\bibitem[31]{published-31} Shuichi Hirahara and Osamu Watanabe. On nonadaptive security reductions of hitting set generators. In Proceedings of 24th International Conference on Randomization and Computation, volume 176 of LIPIcs, pages 15:1–15:14. Schloss Dagstuhl - Leibniz-Zentrum für Informatik, 2020.
\bibitem[32]{published-32} John M. Hitchcock. Limitations of efficient reducibility to the kolmogorov random strings. Comput., 1(1):39–43, 2012.
\bibitem[33]{published-33} John M. Hitchcock and Aduri Pavan. On the NP-completeness of the minimum circuit size problem. In Conference on Foundation of Software Technology and Theoretical Computer Science (FSTTCS), pages 236–245, 2015.
\bibitem[34]{published-34} Rahul Ilango. Approaching MCSP from above and below: Hardness for a conditional variant and ac\textasciicircum{}0[p]. In Proceedings of 11th Innovations in Theoretical Computer Science Conference, volume 151 of LIPIcs, pages 34:1–34:26. Schloss Dagstuhl - Leibniz-Zentrum für Informatik, 2020.
\bibitem[35]{published-35} Rahul Ilango. Constant depth formula and partial function versions of MCSP are hard. In Proceedings of 61st IEEE Annual Symposium on Foundations of Computer Science, pages 424–433, 2020.
\bibitem[36]{published-36} Rahul Ilango, Bruno Loff, and Igor Carboni Oliveira. Np-hardness of circuit minimization for multi-output functions. In Proceedings of 35th Computational Complexity Conference, volume 169 of LIPIcs, pages 22:1–22:36. Schloss Dagstuhl - Leibniz-Zentrum für Informatik, 2020.
\bibitem[37]{published-37} Valentine Kabanets and Jin-yi Cai. Circuit minimization problem. In Symposium on Theory of Computing (STOC), pages 73–79, 2000.
\bibitem[38]{published-38} Adam R. Klivans and Dieter van Melkebeek. Graph nonisomorphism has subexponential size proofs unless the polynomial-time hierarchy collapses. SIAM J. Comput., 31(5):1501–1526, 2002.
\bibitem[39]{published-39} Ker-I Ko. On the complexity of learning minimum time-bounded turing machines. SIAM Journal on Computing, 20(5):962–986, 1991.
\bibitem[40]{published-40} Leonid Levin. Universal search problems. Problems of Information Transmission, 9(3):115–116, 1973.
\bibitem[41]{published-41} William J. Masek. Some NP-complete set covering problems. Unpublished Manuscript, 1979.
\bibitem[42]{published-42} Cody Murray and Ryan Williams. On the (non) NP-hardness of computing circuit complexity. In Conference on Computational Complexity (CCC), pages 365–380, 2015.
\bibitem[43]{published-43} Cody Murray and Ryan Williams. On the (non) np-hardness of computing circuit complexity. Theory Comput., 13(1):1–22, 2017.
\bibitem[44]{published-44} Michael Saks and Rahul Santhanam. Circuit lower bounds from np-hardness of MCSP under turing reductions. In Shubhangi Saraf, editor, Proceedings of 35th Computational Complexity Conference, volume 169 of LIPIcs, pages 26:1–26:13. Schloss Dagstuhl - Leibniz-Zentrum für Informatik, 2020.
\bibitem[45]{published-45} Rahul Santhanam. Pseudorandomness and the minimum circuit size problem. In Thomas Vidick, editor, 11th Innovations in Theoretical Computer Science Conference, ITCS 2020, January 12-14, 2020, Seattle, Washington, USA, volume 151 of LIPIcs, pages 68:1–68:26. Schloss Dagstuhl - Leibniz-Zentrum für Informatik, 2020.
\bibitem[46]{published-46} Stathis Zachos. Probabilistic quantifiers and games. J. Comput. Syst. Sci., 36(3):433–451, 1988. doi:10.1016/0022-0000(88)90037-2.
\end{thebibliography}
\end{document}
